% !TeX root = ../../Book2.tex
% 纲要标题：### 13.2 半单环
% 纲要位置：计划纲要.md，第 1470 行。

\section{半单环}
\label{b2:sec:1302}

半单模可以定义在任意环上；半单环则要求环自身的正则模已经半单。因为每个模都由自由模取商得到，这个看似只涉及一个模的要求，实际上会迫使整个模范畴中的子模都能分裂。它还蕴含强烈的有限性，不能与仅仅“没有非零元素同时零化所有简单模”混淆。

% 纲要要点（供目录核对）：
%
% * 正则模半单的等价条件
% * 所有模的投射性与内射性
% * 半单性不等于半本原性
%

\subsection{正则模半单的等价条件}
\label{b2:subsec:1302:01}

\begin{definition}\label{b2:def:semisimple-ring}
设 \(R\) 为含幺环。若左正则模 \({}_RR\) 是半单模，就称 \(R\) 为\term[bandan-huan]{半单环}[semisimple ring]。零环的正则模为零，按此约定也是半单环。\secref{b2:sec:1303}将证明该定义与右正则模半单等价。
\end{definition}

\begin{theorem}\label{b2:thm:semisimple-ring-characterizations}
设 \(R\) 为含幺环，所论左 \(R\) 模均为幺性模。下列条件等价：
\begin{equivalences}
\item 左正则模 \({}_RR\) 半单。
\item 每个左理想都是 \({}_RR\) 的直和因子。
\item 每个左 \(R\) 模都是半单模。
\item 每条左 \(R\) 模的短正合列都分裂。
\end{equivalences}
若这些条件成立，\({}_RR\) 是有限多个简单左理想的直和，因而具有有限长度，特别是左 Noether 且左 Artin。
\end{theorem}
\begin{proof}
第一、二条的等价是\cref{b2:thm:semisimple-module-characterizations}应用于 \({}_RR\)，因为正则模的子模恰是左理想。假设第一条，要从这个正则模推广到任意左模 \(M\)，可先把它的生成元放进自由模：自由模 \(R^{(I)}\) 是若干份 \({}_RR\) 的直和，由\cref{b2:prop:semisimple-subquotients}仍半单。任意左模 \(M\) 都有满射 \(R^{(I)}\to M\)，例如按 \(M\) 的元素取自由生成元并将其送到对应元素。再用同一命题中商模保持半单的结论，得到第三条。第三条当然包含第一条。

若第三条成立，短正合列 \(0\to A\xrightarrow{i}B\xrightarrow{q}C\to0\) 中，\(i(A)\) 在 \(B\) 内有补模 \(K\)。\(q|_K\) 单射是因为 \(K\cap\ker q=0\)，满射是因为每个 \(b\) 可写成 \(i(a)+k\)。其逆后接 \(K\hookrightarrow B\) 便是截面，所以序列分裂。反过来，将第四条用于每条 \(0\to N\to M\to M/N\to0\)，由\cref{b2:thm:split-exact-sequence}得到每个子模都有补模，故每个 \(M\) 半单。

最后，\(R\) 由 \(1\) 生成，\cref{b2:prop:semisimple-subquotients}中的有限生成判据便使简单分量只有有限多个。具体地，若 \(R=\bigoplus_{i\in I}S_i\)，则 \(1=\sum_{i\in F}s_i\) 只涉及有限集 \(F\subseteq I\)。每个 \(S_i\) 是左理想，所以对任意 \(r\in R\)，有 \(r=r1=\sum_{i\in F}rs_i\in\bigoplus_{i\in F}S_i\)。因此这有限个分量已覆盖全部 \(R\)。逐个加上它们得到有限合成列，故 \({}_RR\) 有限长度。由\cref{b2:thm:finite-length-chain-conditions}，左理想的升链与降链都稳定，得到两项有限性结论。
\end{proof}

在此意义下，“半单 Artin 环”中的左 Artin 条件已由半单性推出。对域 \(k\)，无限直积环 \(R=\prod_{i\ge1}k\) 不会仅因每个因子是域就自动成为半单环：其正则模含有严格下降的左理想链
\[
R\supsetneq\{x:x_1=0\}\supsetneq
\{x:x_1=x_2=0\}\supsetneq\cdots.
\]
每一步都可由仅在相应坐标为 \(1\) 的元素见出严格包含。定理保证半单环必左 Artin，所以这个环不是半单环。它的各个坐标理想虽是简单模，但它们的直和只包含有限支集的元素，不能覆盖整个正则模。

\ProseParagraphBreak
这里讨论的是无限直积环在自身乘法下的正则模，不能据此断言同一个加法群在别的标量作用下也不半单。例如 \(\prod_{i\ge1}\mathbb F_2\) 作为整数模被 \(2\) 零化，因而是一个 \(\mathbb F_2\) 向量空间。选取向量空间基后，它就是一维 \(\mathbb F_2\) 空间的直和，每个分量作为整数模都同构于简单模 \(\mathbb Z/2\mathbb Z\)，故这个整数模半单。这里选取的基不是各个坐标向量组成的集合；后者只生成有限支集的向量。半单性取决于底环及其作用。

\subsection{所有模的投射性与内射性}
\label{b2:subsec:1302:02}

\begin{proposition}\label{b2:prop:semisimple-projective-injective}
设 \(R\) 为含幺环，所论左 \(R\) 模均为幺性模。若 \(R\) 半单，则每个左 \(R\) 模既投射又内射。反过来，若每个左 \(R\) 模都投射，或每个左 \(R\) 模都内射，则 \(R\) 半单。
\end{proposition}
\begin{proof}
设 \(R\) 半单。给定满射 \(q:B\to C\)，其核所确定的短正合列分裂，故有截面 \(s:C\to B\)。对任意 \(f:P\to C\)，映射 \(sf:P\to B\) 是提升，说明任意 \(P\) 投射。给定单射 \(i:A\to B\)，分裂性给出缩回 \(r:B\to A\)；任意 \(f:A\to E\) 都由 \(fr:B\to E\) 延拓，说明任意 \(E\) 内射。

反过来，若所有模都投射，对每条短正合列的商模使用提升性质，将其恒等映射提升为截面，所以所有短正合列都分裂。若所有模都内射，则将子模的恒等映射沿包含延拓，得到缩回，也使每条短正合列分裂。两种情形均由\cref{b2:thm:semisimple-ring-characterizations}推出半单性。
\end{proof}

这里要求“每个左模”都具有性质。给定模 \(M\) 的半单性只保证 \(M\) 自身的子模有补模；投射性却要求对任意满射提升从 \(M\) 出发的同态，内射性要求沿任意单射延拓以 \(M\) 为目标的同态，条件涉及其他模。因而一个给定模半单，不意味着它在原环上投射或内射。例如对素数 \(p\)，简单整数模 \(\mathbb Z/p\mathbb Z\) 不是投射模，因为 \(\mathbb Z\to\mathbb Z/p\mathbb Z\) 没有截面：它的任何截面都会把非零有限阶元素嵌入无挠模 \(\mathbb Z\)。它也不是内射模，因为从 \(p\mathbb Z\) 到它的同态 \(p\mapsto\bar1\) 不能延拓到 \(\mathbb Z\)：任何延拓都会把 \(p\) 送到 \(p\) 倍的某个元素，即零。

\ProseParagraphBreak
半单环上的直和分解通常不唯一到具体子模。例如 \(k^2\) 中 \(ke_1\) 的补空间可取任意 \(k(e_2+ae_1)\)。半单性保证补模存在，结构定理则描述简单类型及其重数；这两者都不要求为每个子模指定唯一补模。

\subsection{半单性与半本原性}
\label{b2:subsec:1302:03}

每个简单模 \(S\) 给出一个双边理想 \(\operatorname{Ann}_R(S)\)，由同时零化 \(S\) 全部元素的标量组成。把这些理想取交，可以检测一个环的所有简单模共同看不见哪些元素。

\ProseParagraphBreak
必须区别零化整个 \(S\) 与零化一个选定向量。对 \(0\ne s\in S\)，满射 \(R\to S\)、\(r\mapsto rs\) 的核 \(L_s=\{\,r\in R\mid rs=0\,\}\) 是极大左理想，而 \(\operatorname{Ann}_R(S)=\bigcap_{s\ne0}L_s\) 要求所有向量同时被零化。即使 \(S\cong R/L_s\)，也不能一般地写成 \(\operatorname{Ann}_R(S)=L_s\)。前面的 \(M_2(k)\) 列模就是例子：对 \(s=(1,0)^{\mathsf t}\)，\(L_s\) 是第一列为零的矩阵组成的非零左理想；但一个矩阵若零化全部列向量，其每一列都为零，故 \(\operatorname{Ann}_{M_2(k)}(k^2)=0\)。

\ProseParagraphBreak
所有简单模的同构类可以用一个集合表示：每个简单模都同构于某个 \(R/L\)，而极大左理想组成 \(R\) 的幂集中的一个集合。从这些商模中每个同构类选一个代表，便得到所需代表族；同构的模具有相同的零化子。

\begin{proposition}\label{b2:prop:simple-annihilator-intersection}
设 \(R\) 为非零含幺环，\(\Sigma\) 为幺性简单左 \(R\) 模的一组同构类代表。有
\[
\bigcap_{S\in\Sigma}\operatorname{Ann}_R(S)
=\bigcap_{L\text{ 为 }R\text{ 的极大左理想}}L.
\]
\end{proposition}
\begin{proof}
若 \(a\) 零化每个简单模，对每个极大左理想 \(L\)，商模 \(R/L\) 简单，故 \(a(1+L)=0\)，即 \(a\in L\)。这证明左边包含于右边。

反过来，设 \(a\) 属于全部极大左理想。固定任意简单模 \(S\)，但让 \(0\ne s\in S\) 遍历其所有非零向量。对每个这样的 \(s\)，满射 \(R\to S\)、\(r\mapsto rs\) 的核 \(L_s\) 都是极大左理想，因而包含 \(a\)，所以 \(as=0\)。这里可能需要随 \(s\) 改换 \(L_s\)，但 \(a\) 属于全部极大左理想的假设对每个 \(s\) 都适用。零向量也被零化，于是 \(a\in\operatorname{Ann}_R(S)\)，得到反向包含。
\end{proof}

若含幺环 \(R\) 的全部幺性简单左模的零化子之交为零，则称 \(R\) 具有\term[banbenyuan-xing]{半本原性}[semiprimitivity]。这就是说，简单模的作用族共同忠实：任何非零 \(a\in R\) 都在某个简单模的某个向量上作用非零；单个简单模的作用则未必忠实。第十四章将把这个交作为 Jacobson 根研究。

\ProseParagraphBreak
半单环一定满足这个条件。若 \(a\) 零化所有简单模，就零化左正则模的每个简单直和分量，因而零化整个 \(R\)；特别是 \(a=a1=0\)。但反过来不成立：\(\mathbb Z\) 的简单模为全部 \(\mathbb Z/p\mathbb Z\)，零化子为 \(p\mathbb Z\)，且
\[
\bigcap_{p\text{ 素数}}p\mathbb Z=0.
\]
若非零整数 \(a\) 属于此交，取 \(|a|+1\) 的任一素因子 \(p\)，它不整除 \(a\)，立即矛盾。但短正合列
\[
0\longrightarrow2\mathbb Z\longrightarrow\mathbb Z
\longrightarrow\mathbb Z/2\mathbb Z\longrightarrow0
\]
不分裂：若 \(2\mathbb Z\) 有补模 \(K\)，则 \(K\cong\mathbb Z/2\mathbb Z\)，而 \(\mathbb Z\) 中没有非零有限阶元素，不可能有这样的 \(K\)。因此 \(\mathbb Z\) 不半单。半本原性保证非零环元素能被某个简单作用检测，却没有保证子模与商模之间的扩张分裂。

\ProseParagraphBreak
零环没有非零简单模。若把空族理想的交按约定取为全环，上述交在零环中仍是零，半单性与半本原性在这个退化情形中均成立。后续结构定理将把零环对应于没有矩阵因子的空乘积，避免把它误列为除环。
